\part{Свойства определенного интеграла}
\section{Свойства определенного интеграла}

\begin{prop} выражаемые равенством: \\
    1. $ \int\limits_a^b f(x) \, dx = -\int\limits_b^a f(x) \, dx $. \\
    2. $ \int\limits_a^b 0 \, dx = 0 $. \\ 
    3. $ \int\limits_a^b 1 \, dx = b - a $. \\ 
    4. $ \text{аддитивность}: \forall f, g \in R[a, b] \Rightarrow f + g \in R[a, b] $. \\
    5. $ \text{однородность}: \forall f \in R[a, b], \forall \alpha \in \mathbb{R} \Rightarrow \alpha f \in R[a, b] $ \text{и} 
    $ \int\limits_a^b \alpha f(x) \, dx = \alpha \int\limits_a^b f(x) \, dx $.
\end{prop}

Доказательство данных свойств следует из определения определенного интеграла через сумму Римана.
(внимательному читателю не составит труда самостоятельно проверить эти свойства).

\begin{prop}
4+5. $ \text{линейность}: \forall \alpha_1, \alpha_2, ..\alpha_k \in \mathbb{R}, \forall f_1, f_2, ..f_k \in R[a, b] \Rightarrow \sum\limits_{i=1}^k \alpha_i f_i \in R[a, b] $. \\
    и $ \int\limits_a^b \sum\limits_{i=1}^k \alpha_i f_i(x) \, dx = \sum\limits_{i=1}^k \alpha_i \int\limits_a^b f_i(x) \, dx $. \\
6. $ f \in R[a,b] \forall [c,d] \subseteq [a,b] \Rightarrow f \in R[c,d] $. Доказательство строится на том, что любое разбиение отрезка $[c,d]$ можно расширить до разбиения отрезка $[a,b]$, а любая интегральная сумма для разбиения отрезка $[c,d]$ будет интегральной суммой для расширенного разбиения отрезка $[a,b]$. \\
7. Аддитивность по промежуткам интегрирования: $ \forall c \in (a, b), \forall f \in R[a, b] \Rightarrow f \in R[c, b] $ и $ \int\limits_a^b f(x) \, dx = \int\limits_a^c f(x) \, dx + \int\limits_c^b f(x) \, dx $.
\end{prop}

\begin{paint}[Геометрический смысл]

    \begin{tikzpicture}[scale=1.2]

    % Оси
    \draw[->] (-0.5,0) -- (5,0) node[right] {$x$};
    \draw[->] (0,-0.5) -- (0,3.5) node[above] {$y$};

    % Точки a и b
    \coordinate (A) at (1,0);
    \coordinate (B) at (4,0);


    % График функции (кривая)
    \draw[thick, red]
        (1,1) 
        .. controls (1.5,2.2) and (2.5,2.3) .. (3,2)
        .. controls (3.5,1.8) and (3.8,2.2) .. (4,2.4);

    % Заливка области
    \fill[blue!20]
        (1,0) -- (1,1)
        .. controls (1.5,2.2) and (2.5,2.3) .. (3,2)
        .. controls (3.5,1.8) and (3.8,2.2) .. (4,2.4)
        -- (4,0) -- cycle;
        
    \fill[pattern=north east lines]
        (2.5,0) -- (2.5,2.16)
        .. controls (3,2) and (3.5,1.8) .. (4,2.4)
        -- (4,0) -- cycle;

    % Надпись f(x)
    \node at (2.5,3) {$f(x)$};

    % Вертикальные линии
    \draw[thick] (1,0) -- (1,1);
    \draw[thick] (4,0) -- (4,2.4);
    \draw[thick] (2.5,0) -- (2.5,2.16);

    % Точка O
    \node at (-0.2,-0.2) {$O$};
    \node at (1,-0.2) {$a$};
    \node at (4,-0.2) {$b$};
    \node at (2.5,-0.2) {$c$};

    \end{tikzpicture}
\end{paint}

\begin{prop}
 8. Свойства интеграла на симметричных промежутках: $ \forall f \in R[-a, a] $ и $ f $ - чётная функция, то $ \int\limits_{-a}^a f(x) \, dx = 2\int\limits_0^a f(x) \, dx $, а если $ f $ - нечётная функция, то $ \int\limits_{-a}^a f(x) \, dx = 0 $.
\end{prop}

\section{Свойства выражаемые неравенствами}

\begin{prop}
1. $ f \in R[a,b] (a < b), \forall x \in [a,b] f(x) \geqslant 0 \Rightarrow \int\limits_a^b f(x) \, dx \geqslant 0 $.
\begin{proof}
    $\sigma = \sum\limits_{i=0}^{n-1} f(\xi_i) \Delta x_i \geqslant 0$ для всех интегральных сумм, а значит $ \int\limits_a^b f(x) \, dx = \lim\limits_{\lambda \to 0} \sigma \geqslant 0 $.
\end{proof}
2. Монотонность: $ f, g \in R[a,b] (a < b), \forall x \in [a,b] f(x) \geqslant g(x) \Rightarrow \int\limits_a^b f(x) \, dx \geqslant \int\limits_a^b g(x) \, dx $.
\begin{proof}
    $ f(x) - g(x) \geqslant 0, \forall x \in [a,b] \Rightarrow \int\limits_a^b (f(x) - g(x)) \, dx \geqslant 0 \Rightarrow \int\limits_a^b f(x) \, dx - \int\limits_a^b g(x) \, dx \geqslant 0 $.
\end{proof}
3. Ограниченность определённого интеграла: $ f \in R[a,b] (a < b), \forall x \in [a,b] m \leqslant f(x) \leqslant M \Rightarrow m(b-a) \leqslant \int\limits_a^b f(x) \, dx \leqslant M(b-a) $.
\begin{proof}
    $ m \leqslant f(x) \leqslant M, \Rightarrow \int\limits_a^b m \, dx \leqslant \int\limits_a^b f(x) \, dx \leqslant \int\limits_a^b M \, dx \Rightarrow m(b-a) \leqslant \int\limits_a^b f(x) \, dx \leqslant M(b-a) $.
\end{proof}
\end{prop}

\begin{paint}[Иллюстрация]

    \begin{tikzpicture}[scale=1.2]

% Оси
\draw[->] (-0.5,0) -- (5,0) node[right] {$x$};
\draw[->] (0,-0.5) -- (0,3.5) node[above] {$y$};

% График функции
\draw[thick, red]
    (1,1)
    .. controls (1.8,2.5) and (2.5,2.6) .. (3,2.4)
    .. controls (3.5,2.2) and (3.8,2.0) .. (4,1.8);

% --- НИЖНЯЯ ЧАСТЬ (двойная штриховка до y=2) ---
\begin{scope}
\clip (0,0) rectangle (5,1.3);

\fill[pattern=north east lines]
    (1,0) -- (1,1)
    .. controls (1.8,2.5) and (2.5,2.6) .. (3,2.4)
    .. controls (3.5,2.2) and (3.8,2.0) .. (4,1.8)
    -- (4,0) -- cycle;

\fill[pattern=north west lines]
    (1,0) -- (1,1)
    .. controls (1.8,2.5) and (2.5,2.6) .. (3,2.4)
    .. controls (3.5,2.2) and (3.8,2.0) .. (4,1.8)
    -- (4,0) -- cycle;

\end{scope}

% --- ВЕРХНЯЯ ЧАСТЬ (одинарная штриховка выше y=2) ---
\begin{scope}
\clip (0,1.3) rectangle (5,5);

\fill[pattern=north east lines]
    (1,0) -- (1,1)
    .. controls (1.8,2.5) and (2.5,2.6) .. (3,2.4)
    .. controls (3.5,2.2) and (3.8,2.0) .. (4,1.8)
    -- (4,0) -- cycle;

\end{scope}

% Линия уровня y=2
\draw[dashed] (0,1.3) -- (4.2,1.3) node[right] {$m$};


\draw[dashed] (0,2.5) -- (4.2,2.5) node[right] {$M$};
% Вертикали
\draw[thick] (1,0) -- (1,1);
\draw[thick] (4,0) -- (4,1.8);

% Подписи
\node at (2.5,3) {$f(x)$};
\node at (-0.2,-0.2) {$O$};
\node at (1,-0.2) {$a$};
\node at (4,-0.2) {$b$};

\end{tikzpicture}
\end{paint}

\begin{prop}
4. $ f \in [a,b] (a < b) \Rightarrow |f| \in R[a,b] \space \text{и} \space |\int\limits_a^b f(x) \, dx| \leqslant \int\limits_a^b |f(x)| \, dx $. \\
Следствие: $ f \in R[a,b] (a < b) \text{и} f(x) \geqslant 0 \forall x \in [a,b]: |f(x)| \leqslant k \Rightarrow | \int\limits_a^b f(x) \, dx | \leqslant k(b-a) $. \\
5. $ f \in R[a,b] (a <b) \text{и} f(x) \geqslant 0 \forall x \in [a,b] \Rightarrow \forall [\alpha, \beta] \subseteq [a,b] \space \int\limits_\alpha^\beta f(x) \, dx \leqslant \int\limits_a^b f(x) \, dx $. \\
\end{prop}

\section{ Теоремы о среднем значении}

\begin{definition}
    Пусть $ f \in R[a,b] (a < b) \Rightarrow \forall x \in [a,b]: m \leqslant f(x) \leqslant M \Rightarrow m(b-a) \leqslant \int\limits_a^b f(x)dx \leqslant M(b-a)$.
    $$m  \leqslant \frac{1}{b-a} \int\limits_a^b f(x)dx \leqslant M.$$
    $$\mu = \frac{1}{b-a} \int\limits_a^bf(x)dx$$ — среднее значение функции на $[a,b]$.
\end{definition}

\begin{theorem}
    о среднем значении.
    $$f \in \mathbb{C}[a,b] \Rightarrow \exists \xi \in [a,b]: \int\limits_a^b f(x)dx = f(\xi)(b-a).$$
    \begin{proof}
        $ f \in C[a,b] \Rightarrow f \in R[a,b]$.
         $$\Downarrow$$
         $\exists \sup f(x) = M = f(x_i)$ и $\inf f(x) = m = f(x_j)$.
         Пусть $ \mu = \frac{1}{b-a} \int\limits_a^b f(x)dx \Rightarrow m \leqslant \mu \leqslant M$. По теореме о промежуточном значении непрерывной функции $\mu$ — значение $f(x)$, т.е. $\exists \, \xi \in [a,b]: f(\xi) = \mu$.
    \end{proof}
\end{theorem}


\begin{task}
Используя свойства определённого интеграла, вычислите 
$$ \int\limits_{-2}^{2} (x^3+2x^2+1)\,dx.$$
\end{task}

\begin{task}
Пусть
$$ 2\leqslant f(x)\leqslant 5,\quad x\in[1,4]. $$
Оцените значение интеграла 
$$ \int\limits_1^4 f(x)\,dx. $$
\end{task}
